% Compilen desde esta carpeta o desde la raíz del repositorio.
\makeatletter
\def\input@path{{../../preamble/}{book2026-2027/preamble/}}
\makeatother
\input{preamble.tex}

\usepackage[spanish]{babel}
\usepackage{lmodern}
\usepackage{titlesec}
\geometry{a4paper,margin=1.6cm}
\linespread{1.10}

\definecolor{HeadingColor}{HTML}{243B53}
\definecolor{InstructionColor}{HTML}{704214}
\definecolor{ConceptColor}{HTML}{205C50}
\definecolor{LinkColor}{HTML}{0047AB}
\hypersetup{allcolors=LinkColor}
\newcommand{\guidelink}[2]{\href{#1}{\textbf{\underline{#2}}}}
\newcommand{\instruction}[1]{\textcolor{InstructionColor}{\textbf{#1}}}
\newcommand{\concept}[1]{\textcolor{ConceptColor}{\textbf{#1}}}
\newcommand{\checkbox}{%
    \begingroup
    \setlength{\fboxrule}{0.4pt}%
    \setlength{\fboxsep}{2pt}%
    \fbox{\rule{0pt}{0.8ex}\rule{0.8ex}{0pt}}%
    \endgroup}

\setlength{\parskip}{4pt}
\setlength{\tabcolsep}{5pt}
\renewcommand{\arraystretch}{1.15}
\setlist{leftmargin=0.65cm,topsep=4pt,partopsep=0pt,parsep=0pt,itemsep=3pt}
\titleformat{\section}{\large\bfseries\color{HeadingColor}}{}{0pt}{}
\titlespacing*{\section}{0pt}{10pt}{3pt}
\newcolumntype{Y}{>{\raggedright\arraybackslash}X}

\hypersetup{pdftitle={C1--C2. Generen alternativas y elijan una solución},pdfsubject={Emprendimiento e Innovación, Grado 10, Periodo 2}}

\begin{document}

\begin{center}
{\Large\bfseries\color{HeadingColor} C1--C2. Generen alternativas y elijan una solución\par}
\vspace{3pt}
Emprendimiento e Innovación \enspace|\enspace Grado 10 \enspace|\enspace Periodo 2\\
Guía 1 de 5
\end{center}

\section*{Objetivo y punto de partida}
\concept{C1.} Propone alternativas de solución variadas frente al reto de innovación definido.\\
\concept{C2.} Selecciona una solución según valor para el usuario, viabilidad y coherencia con el problema.

Recuperen D1 (P1, C7), los hallazgos H sobre el usuario (P1, C5), las ideas I (P1, C9) y las condiciones o preguntas pendientes (P1, C10). Partan de la necesidad, el usuario y el contexto definidos en D1.

\section*{Documento acumulativo e integridad de la evidencia}
Esta guía es de consulta. Continúen el \textbf{mismo documento del proyecto en LaTeX} en \guidelink{https://www.overleaf.com/}{Overleaf} o en su editor habitual. Actualicen la fecha y entreguen el PDF completo. Usen \textbf{P1} para el Periodo 1 y \textbf{P2} para el Periodo 2 al citar criterios. Enlacen los registros por código o referencia interna y conserven sus fuentes en la sección de origen. Por ejemplo, si la etiqueta de D1 es \verb|sec:desafio|, usen \verb|\hyperref[sec:desafio]{D1}|.

Distingan \concept{evidencia} (lo observado o expresado), \concept{conclusión} (lo que infieren), \concept{supuesto} (lo pendiente de verificar) e \concept{idea} (una propuesta). Registren fuentes y fechas reales. Identifiquen toda simulación autorizada por el docente. \instruction{Si usan IA, limítenla al código LaTeX, al formato y a organizar contenido propio.} Evidencias, conclusiones, ideas, decisiones y aportes deben provenir de su trabajo.

\section*{C1 de P2. Generen alternativas con mecanismos distintos}
\begin{enumerate}[series=pasos]
\item \textbf{Recuperen el reto y sus fuentes.} Localicen D1 y los H pertinentes. Revisen qué sabe el equipo sobre lo que el usuario necesita lograr y qué sigue siendo un supuesto. Citen los códigos al describir qué parte del reto aborda cada alternativa.
\item \textbf{Amplíen las ideas.} Retomen, combinen o desarrollen ideas I y planteen al menos tres alternativas. Para cada una, expliquen qué haría el usuario o el servicio y cómo esa acción produciría la mejora esperada. Si dos ideas funcionan igual, busquen otra forma de resolver la dificultad.
\item \textbf{Registren las alternativas en la matriz de la página siguiente.} Denles códigos ALT1, ALT2, etc. Enlacen las ideas I de origen cuando corresponda y el hallazgo H que orienta cada propuesta.
\end{enumerate}

\textbf{Ejemplo ficticio.} Si el reto fuera reducir la espera en una cafetería escolar, podrían distribuir las llegadas por turnos, preparar pedidos por adelantado o separar la entrega del pago. Cada opción interviene en un momento o actividad distinta del servicio.

\concept{Evidencia de C1.} Tres o más alternativas con mecanismos distinguibles, vinculadas a D1 y a los registros que orientaron su creación.

\newpage
\section*{C2 de P2. Comparen y justifiquen la elección}
\begin{enumerate}[resume=pasos]
\item \textbf{Comparen con las mismas preguntas.} Para valorar el aporte al usuario, pregúntense qué mejora y qué H respalda esa expectativa. Para valorar la viabilidad, revisen tiempo, materiales, acceso y permisos. Para valorar la coherencia, expliquen qué parte de D1 aborda el mecanismo. Usen \textbf{SD (sin datos)} cuando falte respaldo y anoten qué necesitan saber.
\end{enumerate}

{\small
\begin{tabularx}{\linewidth}{|Y|Y|Y|Y|}
\hline
\textbf{Alternativa y mecanismo} & \textbf{Valor para el usuario} & \textbf{Viabilidad} & \textbf{Coherencia con D1} \\
\hline
Código ALT, cómo funciona y origen I, si aplica. & Beneficio esperado y hallazgo H que lo sustenta. & Recursos disponibles, barreras o condiciones pendientes. & Parte del reto que atiende y cómo la atiende. \\
\hline
\end{tabularx}\par
}

\begin{enumerate}[resume=pasos]
\item \textbf{Elijan y expliquen qué inclinó la decisión.} Nombren la alternativa y contrástenla con al menos otra. Expliquen qué evidencia del usuario y qué condición de viabilidad pesaron más. Indiquen qué supuesto podría cambiar la elección y cómo lo comprobarían. Si ese supuesto es decisivo, marquen la elección como provisional.
\item \textbf{Abran la bitácora del equipo.} Desde hoy y al cierre de cada sesión, registren fecha, quién aportó qué y a qué registro se vincula, compromisos con responsable y plazo, y su estado. Cuando surja un desacuerdo, registren las opciones discutidas y cómo llegaron a un acuerdo o qué quedó pendiente. Pueden usar el código B1, B2, etc., para cada entrada.
\end{enumerate}

\concept{Evidencia de C2.} Una matriz comparativa y una decisión argumentada con respaldo identificable, condiciones de viabilidad y supuestos pendientes.

\section*{Resultado y continuidad}
Agreguen la matriz, la decisión y la primera entrada de bitácora al documento acumulativo. En P2, C3--C4 recuperarán la alternativa elegida y el supuesto prioritario para decidir qué función o interacción representar y prototipar. En P2, C9 volverán a las entradas de la bitácora.

\section*{Verificación final}
\begin{itemize}[label=\checkbox]
\item Las alternativas proponen mecanismos distintos para el mismo D1.
\item Las comparamos por valor para el usuario, viabilidad y coherencia con el reto.
\item La elección contrasta alternativas, cita evidencias y distingue los supuestos pendientes.
\item Abrimos la bitácora y compilamos el PDF acumulativo con sus referencias.
\end{itemize}

\end{document}
